\documentclass[12pt]{article}
\usepackage{amsmath,amssymb}

\begin{document}
\section*{{2018 AIME II Problems/Problem 9}}
Source: AoPS Wiki

\subsection*{Solution}
We represent octagon $ABCDEFGH$ in the coordinate plane with the upper left corner of the rectangle being the origin. Then it follows that $A=(0,-6), B=(8, 0), C=(19,0), D=(27, -6), E=(27, -17), F=(19, -23), G=(8, -23), H=(0, -17), J=(0, -\frac{23}{2})$ . Recall that the centroid is $\frac{1}{3}$ way up each median in the triangle. Thus, we can find the centroids easily by finding the midpoint of the side opposite of point $J$ . Furthermore, we can take advantage of the reflective symmetry across the line parallel to $BC$ going through $J$ by dealing with less coordinates and ommiting the $\frac{1}{2}$ in the shoelace formula. By doing some basic algebra, we find that the coordinates of the centroids of $\bigtriangleup JAB, \bigtriangleup JBC, \bigtriangleup JCD, \bigtriangleup JDE$ are $\left(\frac{8}{3}, -\frac{35}{6}\right), \left(9, -\frac{23}{6}\right), \left(\frac{46}{3}, -\frac{35}{6}\right),$ and $\left(18, -\frac{23}{2}\right)$ , respectively. We'll have to throw in the projection of the centroid of $\bigtriangleup JAB$ to the line of reflection to apply shoelace, and that point is $\left( \frac{8}{3}, -\frac{23}{2}\right)$ Finally, applying Shoelace, we get: $\left|\left(\frac{8}{3}\cdot (-\frac{23}{6})+9\cdot (-\frac{35}{6})+\frac{46}{3}\cdot (-\frac{23}{2})+18\cdot (\frac{-23}{2})+\frac{8}{3}\cdot (-\frac{35}{6})\right) - \left((-\frac{35}{6}\cdot 9) +\\(-\frac{23}{6}\cdot \frac{46}{3})+ (-\frac{35}{6}\cdot 18)+(\frac{-23}{2}\cdot \frac{8}{3})+(-\frac{23}{2}\cdot \frac{8}{3})\right)\right|$ $=\left|\left(-\frac{92}{9}-\frac{105}{2}-\frac{529}{3}-207-\frac{140}{9}\right)-\left(-\frac{105}{2}-\frac{529}{9}-105-\frac{92}{3}-\frac{92}{3}\right)\right|$ $=\left|-\frac{232}{9}-\frac{1373}{6}-207+\frac{529}{9}+\frac{184}{3}+105+\frac{105}{2}\right|$ $=\left|\frac{297}{9}-\frac{690}{6}-102\right|=\left| 33-115-102\right|=\left|-184\right|=\boxed{184}$ Solution by ktong note: for a slightly simpler calculation, notice that the heptagon can be divided into two trapezoids of equal area and a small triangle.

\end{document}
